\documentclass[11pt,a4paper]{article}

\usepackage[margin=0.55in]{geometry}
\usepackage[T1]{fontenc}
\usepackage[utf8]{inputenc}
\usepackage{lmodern}
\usepackage{microtype}
\usepackage{tabularx}
\usepackage{enumitem}
\usepackage[hidelinks]{hyperref}
\usepackage{array}

\pagestyle{empty}
\linespread{1.01}
\setlength{\parindent}{0pt}
\setlength{\tabcolsep}{0pt}
\setlist[itemize]{
  leftmargin=1.15em,
  itemsep=2.0pt,
  topsep=2.8pt,
  parsep=0pt,
  partopsep=0pt
}

% Education entry
\newcommand{\eduentry}[4]{%
  \begin{tabularx}{\linewidth}{@{}X>{\raggedleft\arraybackslash}p{0.27\linewidth}@{}}
    \textbf{#1} & \textbf{#2} \\
    \textit{#3} & \textit{#4}
  \end{tabularx}
}

% Research experience entry
% First row: project / experience name | institution
% Second row: position | date + location
\newcommand{\researchentry}[5]{%
  \begin{tabularx}{\linewidth}{@{}
    >{\raggedright\arraybackslash\hyphenpenalty=10000\exhyphenpenalty=10000}X
    r@{}}
    \textbf{#1} & \textbf{#2}
  \end{tabularx}\\[-1pt]
  \begin{tabularx}{\linewidth}{@{}Xr@{}}
    \textit{#3} & \textit{\mbox{#4 \;|\; #5}}
  \end{tabularx}
}

% Section separator
\newcommand{\cvsection}[1]{%
  \vspace{5pt}
  \hrule height 0.45pt
  \vspace{3pt}
  {\large\bfseries #1}\par
  \vspace{3pt}
}

% Sub-project title
\newcommand{\projecttitle}[1]{%
  \vspace{1pt}
  \textbf{#1}\par
  \vspace{-1pt}
}

\begin{document}

% ============================================================
% Header
% ============================================================
\begin{center}
  {\LARGE\bfseries Ximing Wang}\\[-1pt]
  \href{tel:+8618661675002}{(+86) 186 6167 5002}
  \;|\;
  \href{mailto:bs_wxm@mail.scut.edu.cn}{bs\_wxm@mail.scut.edu.cn}
  \;|\;
  \href{mailto:vexpaer@gmail.com}{vexpaer@gmail.com}
\end{center}

\vspace{-7pt}

% ============================================================
% Education
% ============================================================
\cvsection{Education}

\eduentry
{South China University of Technology}
{Sep. 2023 -- Present}
{B.Eng. in Biomedical Engineering; GPA: 3.77/4.00; Average: 87.79/100}
{Guangzhou, China}

\vspace{2pt}
\textbf{Selected Coursework:}
Signals and Systems, Calculus, Linear Algebra,
Probability and Mathematical Statistics, Physiology,
Anatomy, Biochemistry, Molecular Biology, Biomechanics

% ============================================================
% Academic Paper
% ============================================================
\cvsection{Academic Paper}

\textbf{An EEG Dataset for Motion Sickness Assessment in Electric Vehicle
Passengers under Normal and Accelerated Driving}.
Manuscript under revision at \textit{Scientific Data}.\\[-1pt]
Dataset:
\href{https://doi.org/10.57760/sciencedb.32147}{doi.org/10.57760/sciencedb.32147}

% ============================================================
% Research Interests
% ============================================================
\cvsection{Research Interests}

EEG-based neuroengineering;
biomedical signal processing;
brain-computer interfaces;
machine learning for real-time physiological monitoring

% ============================================================
% Research Experience
% ============================================================
\cvsection{Research Experience}

\researchentry
{Real-World EEG Study of Motion Sickness}
{South China University of Technology}
{Undergraduate Researcher}
{Mar. 2025 -- Sep. 2025}
{Guangzhou, China}

\begin{itemize}
  \item Designed and conducted real-vehicle experiments under high-
  and low-acceleration driving conditions, collecting synchronized
  32-channel EEG, ECG, and EMG signals from 15 participants.
  \item Developed an end-to-end physiological-signal analysis workflow
  covering EEG preprocessing, ICA-based artifact removal, feature
  extraction, and machine-learning modeling.
  \item Developed an online motion-sickness monitoring system using
  synchronized EEG, ECG, and EMG signals.
\end{itemize}

\vspace{7pt}

\researchentry
{Driver Fatigue Analysis and OpenBCI Software Development}
{Tsinghua University}
{Research Assistant Intern}
{Mar. 2026 -- Aug. 2026}
{Beijing, China}

\projecttitle{Driver Fatigue two-channel-EEG Analysis}
\begin{itemize}
  \item Analyzed approximately 500 hours of dual-channel EEG,
  heart-rate, and KSS data from 165 participants collected during
  daytime and nighttime driving sessions.
  \item Built a preprocessing and quality-control pipeline for
  EEG/KSS alignment, artifact screening, segmentation, and continuous
  fatigue-label construction.
  \item Developed fatigue-estimation models using time-frequency
  representations, channel-wise signal decomposition, and contrastive
  representation learning for robust continuous fatigue estimation.
\end{itemize}

\projecttitle{OpenBCI Embedded Firmware and Host Software}
\begin{itemize}
  \item Developed embedded firmware from scratch for an OpenBCI
  acquisition board, implementing biosignal acquisition, board
  configuration, data buffering, device control, and real-time
  data streaming.
  \item Built a Python desktop application for device communication,
  real-time waveform visualization, acquisition control, signal
  recording, and data export.
  \item Developed a companion mobile application for device connection,
  acquisition control, and real-time signal monitoring, completing
  the software workflow from embedded acquisition to desktop and
  mobile interfaces.
\end{itemize}

% ============================================================
% Awards
% ============================================================
\cvsection{Awards}

\textbf{Second Prize},
10th National Biomedical Engineering Innovation and Design Competition,
2025

% ============================================================
% Skills
% ============================================================
\cvsection{Skills}

\textbf{Programming and Tools:}
Python, MATLAB, R, LaTeX, Adobe Illustrator, Photoshop

\textbf{Experimental and Analytical:}
EEG, ECG, EMG, and fNIRS acquisition;
embedded firmware development;
signal preprocessing;
ICA;
feature extraction;
machine-learning modeling

\textbf{Languages:}
Chinese Mandarin (native);
English (IELTS: 7.0/9.0; CET-6: 604/710)

\end{document}
